% !TEX root = ../main.tex
\chapter{Optical Geometry I}\label{ch:c16}
\begin{paracol}{2}

\begin{sources}
\begin{tabularx}{\linewidth}{@{}l l L@{}}
\toprule
\textbf{Source} & \textbf{Length} & \textbf{Content}\\
\midrule
\seg{16}{1}{Lecture 16: Optical Geometry I} & 1:14:06 & Plan of the two lectures (lensing as deflection in a three-dimensional space); Fermat's principle in classical optics and why it fails for null curves; stationary and static spacetimes (timelike Killing field, the general line element); null curves and null geodesics; the variation of the null condition and Fermat's principle in static spacetimes (the arrival time is stationary); the optical metric of static spacetimes, the Randers--Finsler metric of stationary ones; the optical geometry of Schwarzschild, its isometric embedding in $\mathbb R^3$ (the throat at $r=3M$); the Gaussian curvature of a surface and the lemma $K=-f''/f$ with proof, the sphere as a check; the Gaussian curvature of the Schwarzschild optical geometry, negative everywhere, and the question of how light rays can reconverge. Student questions throughout.\\
\bottomrule
\end{tabularx}

\smallskip
\textbf{Transcript} \file{materials/transcripts/C16-V01.txt} (YouTube captions, lecture given on Ash Wednesday, 2015, as said in class). \textbf{Lecturer: Marcus C.\ Werner, not Schuller} (satellite lecture of the school; the chapter continues in chapter~\ref{ch:c17}, after the break of this lecture).
\end{sources}

\begin{coursenote}[Lecturer]
Lecturer: Marcus C.\ Werner, not Schuller. The notation of this guest lecture differs from Schuller's in places (signature $(-,+,+,+)$ instead of Schuller's $(+,-,-,-)$, Greek and Latin indices for spacetime and space, a different sign convention for the Riemann tensor in the proof of \cref{lem:c16-gauss}): it is kept as given in class.
\end{coursenote}

\begin{reading}[Where this chapter sits in the other notes]
A related paper by Schuller and Werner is arXiv:1707.01261, ``Etherington's distance duality with birefringence'' (its content was not checked for this chapter); the references of the lectures are Gibbons and Werner for the Gauss--Bonnet approach to lensing (as said in class, see chapter~\ref{ch:c17}). In this book: Killing vector fields and symmetries are chapter~\ref{ch:c11}, the relativistic spacetime and its null curves chapter~\ref{ch:c13}, the Einstein equation that Schwarzschild solves chapter~\ref{ch:c15}, Schwarzschild itself chapter~\ref{ch:c22} (not written yet). \emph{Theoretische Mechanik} (Schuller's classical mechanics notes): the variational principles and the Fermat-type arguments [chapter not checked]. The curvature chapter~\ref{ch:c08} is not written yet; the Riemann tensor is used here as in chapter~\ref{ch:c07}.
\end{reading}

% ==================================================================
\section{Plan of the lectures and Fermat's principle in optics}\label{sec:c16-plan}

``Good morning everyone, it is Ash Wednesday today and as you can see it is rather bleak'' \ts{16}{1}{0:00}. Yesterday one started talking about lensing (the satellite lecture before); today the \emph{theory behind lensing} is discussed, from the point of view of general relativity. When one talks about the deflection of light one really means deflection \emph{in space}: so one has to find a way of going from the four-dimensional spacetime down to a three-dimensional space in which deflection can be discussed, and that is the subject of \textbf{optical geometry} \ts{16}{1}{0:20}\ts{16}{1}{0:41}.

The plan of the first lecture:
\begin{itemize}
  \item Fermat's principle and the form it takes in general relativity;
  \item optical geometry; to do calculations one specialises to \emph{static} spacetimes, whose optical geometries are \emph{Riemannian};
  \item within this framework lensing can be discussed, and one sees that it is a \emph{global} effect: the topology of the surface that one looks at determines whether one can have multiple images, what the deflection angle is and so on; the technique for global properties is the \textbf{Gauss--Bonnet theorem}, the final topic of the second lecture \ts{16}{1}{1:01}\ts{16}{1}{1:21}\ts{16}{1}{1:43}\ts{16}{1}{2:03}.
\end{itemize}

\paragraph{Fermat's principle in classical optics.} There is an earlier form, going back to antiquity (Heron): light always takes the \emph{shortest path}. This is true only in a homogeneous optical medium; more generally Fermat observed that light always takes the path that is shortest \emph{in time}. If one integrates along possible paths, the time to traverse a length element $\mathrm dl$ is $\mathrm dl$ divided by the speed, and the speed is the speed of light $c$ divided by the refractive index $n$ (which is $1$ in vacuum or a number greater than $1$, and may depend on space):
\[
  T[\text{path}]=\int\frac{n}{c}\,\mathrm dl.
\]
Looking for stationary curves of this, one gets the actual light rays \ts{16}{1}{2:30}\ts{16}{1}{2:50}\ts{16}{1}{3:11}\ts{16}{1}{3:32}.
Of course this formulation does \emph{not} work in general relativity, because the length in spacetime of a null curve is zero. One must find a way of rewriting it \ts{16}{1}{3:32}\ts{16}{1}{3:52}.

\section{Stationary and static spacetimes}\label{sec:c16-static}

Setting up (``I am not sure if your conversation is basing from the lectures in detail, so I set up''): a spacetime is a four-dimensional differentiable manifold with a Lorentzian metric, of signature $(-,+,+,+)$ here, worked with in local charts, with Greek indices $\mu=0,\dots,3$ for spacetime coordinates and Latin indices $a,b,\dots$ for spatial ones \ts{16}{1}{3:52}\ts{16}{1}{4:13}\ts{16}{1}{4:38}\ts{16}{1}{4:59}.

\begin{definition}[New concepts: stationary spacetime]\label{def:c16-stat}
A spacetime is called \textbf{stationary} if it has a \emph{timelike Killing vector field} $T$: $T$ is timelike (so $g(T,T)<0$ in the signature used), and $T$ is Killing, i.e.\ corresponds to a symmetry of the metric, the Lie derivative of the metric along $T$ vanishes,
\[
  \mathcal L_Tg=0.
\]
\end{definition}
\ts{16}{1}{4:59}\ts{16}{1}{5:19}\ts{16}{1}{5:41}\ts{16}{1}{6:02}
(Killing fields and Lie derivatives were seen in the lectures and tutorials: chapter~\ref{ch:c11}.) What does this condition mean? In components $(\mathcal L_Tg)_{\mu\nu}$ has two terms (``if you look at the derivative of $T$ you get two terms like that'') from the derivatives of $T$ and one with $T^\lambda\partial_\lambda g_{\mu\nu}$. Now choose coordinates adapted to the Killing field: $T=\partial/\partial t$, a vector field with components $(1,0,0,0)$. Then the terms with derivatives of $T$ drop out, and what is left is the derivative of the metric components with respect to $t$: $\partial_tg_{\mu\nu}=0$ for $T$ to be Killing. (One does not really have to use Killing's equation in this case.) So \emph{we can find a coordinate system in which the metric components do not depend on time} \ts{16}{1}{6:02}\ts{16}{1}{6:27}\ts{16}{1}{6:48}\ts{16}{1}{7:10}\ts{16}{1}{7:31}\ts{16}{1}{7:52}.

The most general stationary spacetime can then be written as a line element
\begin{equation}\label{eq:c16-stat}
  \mathrm ds^2=-V^2\bigl(\mathrm dt+\omega_a\,\mathrm dx^a\bigr)^2+h_{ab}\,\mathrm dx^a\mathrm dx^b,
\end{equation}
with a Riemannian metric $h_{ab}$ on three-dimensional space (signature $(+,+,+)$), a mixing term between time and space, and functions $V$, $\omega_a$ and $h_{ab}$ that do \emph{not} depend on $t$ [the symbols $V,\omega,h$ and the exact placement of the squares are taken from the captions and the later use, not from a readable board: $g_{00}=-V^2$, $g_{0a}=-V^2\omega_a$ is one reading, ``all the possible metric coefficients are covered with this''] \ts{16}{1}{7:52}\ts{16}{1}{8:18}\ts{16}{1}{8:43}\ts{16}{1}{9:05}\ts{16}{1}{9:26}\ts{16}{1}{10:05}.

\begin{definition}[New concepts: static spacetime]\label{def:c16-static}
A spacetime is \textbf{static} if it has a timelike Killing vector field $T$ that is moreover \emph{hypersurface orthogonal}. Defining this properly needs the condition of Frobenius's theorem, ``a little bit fiddly; I do not need this for what we are going to do later''; for our purposes it means that \emph{in a particular coordinate system there are no mixing terms}: $g_{0a}=0$, i.e.\ $\omega_a=0$ in \eqref{eq:c16-stat}, and the line element is
\[
  \mathrm ds^2=-V^2\,\mathrm dt^2+h_{ab}\,\mathrm dx^a\mathrm dx^b.
\]
\end{definition}
\ts{16}{1}{10:33}\ts{16}{1}{10:54}\ts{16}{1}{11:16}\ts{16}{1}{11:37}
The $\omega_a$ have a nice symmetric interpretation: they are the spatial part of what is called the \emph{twist vector} of the spacetime \ts{16}{1}{11:37}\ts{16}{1}{11:58}.
\begin{coursenote}[Stationary versus static is a property of the whole spacetime]
The property ``static'' or ``stationary'' is a property of the \emph{entire spacetime}: only certain spacetimes are static or stationary. In contrast, a ``stationary observer'' (one fixed with respect to the coordinates, rather than freely falling, in the space of a black hole) can be defined in various ways. A student asks about the limit of staticness: Schwarzschild is static only outside the horizon $r=2M$ (everything done here is outside of $2M$); the rotating black hole (Kerr) is stationary but not static \ts{16}{1}{37:31}\ts{16}{1}{37:52}\ts{16}{1}{38:13}\ts{16}{1}{51:08}\ts{16}{1}{51:28}.
\end{coursenote}

\section{Null curves and Fermat's principle in static spacetimes}\label{sec:c16-fermat}

Recall \textbf{null curves}: maps $\gamma$ from an interval of the real line into the manifold, in a parametrisation $x^\mu(l)$ (``$l$ for length if you like''), whose tangent vector $\dot x^\mu=\mathrm dx^\mu/\mathrm dl$ has length zero in the spacetime metric,
\[
  g_{\mu\nu}\,\dot x^\mu\dot x^\nu=0.
\]
Among all null curves, certain ones are \textbf{null geodesics}: those that in addition satisfy the geodesic equation $\nabla_{\dot\gamma}\dot\gamma=0$ (``autoparallels'') \ts{16}{1}{11:58}\ts{16}{1}{12:19}\ts{16}{1}{12:39}\ts{16}{1}{13:01}\ts{16}{1}{13:21}\ts{16}{1}{13:41}.

\paragraph{The variation.} ``Now I would like to move on to Fermat's principle in general relativity'' \ts{16}{1}{14:03}. Consider the integral $\int_0^1L\,\mathrm dl$ with $L=g_{\mu\nu}\dot x^\mu\dot x^\nu$ along null curves, the parameter running from $0$ to $1$ (``I would also like to fix the interval''): it is zero for all null curves. Take the variation $\delta x^\mu$ of the curve. As in the derivation of the geodesic equation (``you have probably seen this'') one gets the integral of the Euler--Lagrange expression times the variation, plus the boundary terms at the start and end point:
\[
  0=\delta\int_0^1L\,\mathrm dl=\int_0^1\Bigl(\text{E.L.}\Bigr)_\mu\delta x^\mu\,\mathrm dl+\Bigl[\frac{\partial L}{\partial\dot x^\mu}\,\delta x^\mu\Bigr]_0^1.
\]
``Writing this down explicitly is slightly tedious, I am not going to do it on the blackboard; if you have problems doing this we can do it later.'' The result: since $\partial L/\partial\dot x^\mu=2g_{\mu\nu}\dot x^\nu$ (two terms by the symmetry of $g$), the boundary term is $2g(\dot\gamma,\delta x)$, and the Euler--Lagrange expression is, up to an overall factor $2$ and a minus sign, $g(\delta x,\nabla_{\dot\gamma}\dot\gamma)$, the geodesic equation. Hence
\begin{equation}\label{eq:c16-star}
  g\bigl(\dot\gamma,\delta x\bigr)\Big|_0^1=\int_0^1g\bigl(\delta x,\nabla_{\dot\gamma}\dot\gamma\bigr)\,\mathrm dl\tag{$\ast$}
\end{equation}
(up to the factor and the sign as written on the board [?]) \ts{16}{1}{14:03}\ts{16}{1}{14:33}\ts{16}{1}{14:53}\ts{16}{1}{15:15}\ts{16}{1}{15:43}\ts{16}{1}{16:03}\ts{16}{1}{16:26}\ts{16}{1}{16:47}\ts{16}{1}{17:08}\ts{16}{1}{17:31}\ts{16}{1}{17:52}.

\paragraph{Case 1: both end points fixed.} Suppose all curves have no variation at the beginning of the interval and also at the end point, $\delta x=0$ at $l=0$ and $l=1$. The boundary term vanishes and one has $\int_0^1g(\delta x,\nabla_{\dot\gamma}\dot\gamma)\,\mathrm dl=0$ for all variations $\delta x(l)$, which implies $\nabla_{\dot\gamma}\dot\gamma=0$ (``the metric is non-degenerate, yes, that is the point''): we recover the geodesic equation, which is of course something one knows \ts{16}{1}{18:13}\ts{16}{1}{18:36}\ts{16}{1}{18:56}\ts{16}{1}{19:17}\ts{16}{1}{19:43}\ts{16}{1}{21:10}\ts{16}{1}{21:30}.

\paragraph{Case 2: a static spacetime and the variation at the end in the time direction.} Now suppose one is in a static spacetime, with a coordinate system in which time and space split; the variation at the start point is zero, and the variation \emph{at the end point is only in the time direction}: $\delta x|_1=\delta t\,T$. The picture: all curves start at the same point, but at the parameter value $1$ the variation is just along the Killing vector $T=\partial_t$ \ts{16}{1}{19:43}\ts{16}{1}{20:04}\ts{16}{1}{20:25}\ts{16}{1}{20:48}\ts{16}{1}{21:51}\ts{16}{1}{22:12}. Then \eqref{eq:c16-star} reads, in components,
\[
  g_{\mu\nu}\,\delta x^\mu\,\dot\gamma^\nu\Big|_{l=1}-g_{\mu\nu}\,\delta x^\mu\,\dot\gamma^\nu\Big|_{l=0}=g_{\mu\nu}\,\delta t\,T^\mu\,\dot\gamma^\nu\Big|_{l=1}=\int_0^1g(\delta x,\nabla_{\dot\gamma}\dot\gamma)\,\mathrm dl,
\]
the term at $0$ being zero by assumption. Now \emph{the spacetime is static}: there are no mixed terms, $g_{0a}=0$, so of the sum only one term survives, $g_{00}\,\delta t\,\dot\gamma^0|_1$. In a static spacetime one can assume $g_{00}\neq0$ and that the time component of the tangent vector is non-zero (the curve is not purely spatial), and this is the crucial observation: one is allowed to divide by these two terms,
\begin{equation}\label{eq:c16-deltat}
  \delta t=\frac{1}{g_{00}\,\dot\gamma^0\big|_{l=1}}\int_0^1g\bigl(\delta x,\nabla_{\dot\gamma}\dot\gamma\bigr)\,\mathrm dl.
\end{equation}
This is the outcome of the variation for an arbitrary null curve: \emph{the variation of the arrival time}, in the Killing direction, at the end point. It holds in static spacetimes and not necessarily in others \ts{16}{1}{22:40}\ts{16}{1}{23:02}\ts{16}{1}{23:24}\ts{16}{1}{23:49}\ts{16}{1}{24:10}\ts{16}{1}{24:41}\ts{16}{1}{25:03}\ts{16}{1}{25:26}\ts{16}{1}{25:47}.

\begin{theorem}[Fermat's principle in static spacetimes]\label{thm:c16-fermat}
In a static spacetime, along a null geodesic the variation \eqref{eq:c16-deltat} of the arrival time $t$ at the end point vanishes: $\delta t=0$ (the arrival time is \emph{stationary}) among the null curves emitted at the same point. (The converse is not said in class; it is added in the proof.)
\end{theorem}
\begin{proof}
If the variation is about a null geodesic, then $\nabla_{\dot\gamma}\dot\gamma=0$ and the right-hand side of \eqref{eq:c16-deltat} vanishes: $\delta t=0$ at the end point. This is the ``innocuous result''; (not said in class) conversely if $\delta t=0$ for all variations the integral in \eqref{eq:c16-deltat} vanishes for all $\delta x$, which gives the geodesic equation as in case~1. (Werner says: ``for static spacetimes you have basically seen it, we have gone over the argument''; for stationary spacetimes the proof is ``actually quite complicated'' and is not given: the theorem holds there too.) \ts{16}{1}{26:08}\ts{16}{1}{26:29}\ts{16}{1}{26:51}\ts{16}{1}{27:13}\ts{16}{1}{28:16}\ts{16}{1}{28:37}\ts{16}{1}{29:04}\ts{16}{1}{29:29}
\end{proof}
This looks rather like Fermat's principle seen before; ``so what exactly is different? Does anyone want to comment?'' The answer given by a participant: in ordinary optics what is stationary is the \emph{total travel time} along the whole curve; in general relativity the length of a null curve is zero, so it is not the total travel time but the \textbf{arrival time} that is stationary, ``that is a local property'': it makes sense in general relativity, because one is not looking at the whole of the null curve but just at the arrival time \ts{16}{1}{27:13}\ts{16}{1}{27:35}\ts{16}{1}{27:56}\ts{16}{1}{28:16}.
\begin{coursenote}[Student question: do two rays arrive at the same time?]
A student asks (the question is partly lost in the captions) whether two light rays from the same source arrive at the same time: ``not necessarily, it depends on whether they were emitted at the same time; I can show you an example later'' (the time delay between lensed images appears at the end of the next lecture, chapter~\ref{ch:c17}) \ts{16}{1}{29:29}\ts{16}{1}{30:18}.
\end{coursenote}

\section{The optical metric}\label{sec:c16-optical}

``Now let us look at what this means for our problem of defining spatial light rays in generality.'' By Fermat's principle light rays are the curves satisfying the stationarity condition; what is $t$ for stationary spacetimes? Take the general line element \eqref{eq:c16-stat}, impose the null condition $\mathrm ds^2=0$ and solve for $\mathrm dt$:
\[
  \mathrm dt=-\omega_a\,\mathrm dx^a\pm\frac1V\sqrt{h_{ab}\,\mathrm dx^a\mathrm dx^b}.
\]
(The sign is that of a future-directed ray; the formula is written from the structure of \eqref{eq:c16-stat}, not read from the board [?].) This is the time interval of the line element; to determine the light rays one integrates it, keeping the emission time fixed, and varies: ``light rays are of course just the geodesics of this line element'' \ts{16}{1}{30:39}\ts{16}{1}{31:01}\ts{16}{1}{31:25}\ts{16}{1}{31:52}\ts{16}{1}{32:12}.

This looks unfamiliar: besides the Riemannian square root there is a one-form term $\omega_a\mathrm dx^a$, so $\mathrm dt$ is not given by a Riemannian line element. The geometry of the spatial light rays in stationary spacetimes is not Riemannian at all, ``rather interesting'': it is a more general \textbf{Finsler geometry}, of the particular type called \textbf{Randers} (Randers introduced a line element of this form in the 1940s, as said in class, ``to introduce a preferred direction of time'' in a spacetime context; not checked). This is what describes light rays in the Kerr spacetime (rotating black holes, if one gets to them later in the course); the rotation gives a preferred direction, rays going one way or the other way have different paths, and this is encoded in $\omega_a$, which depends on the angular momentum of the black hole. ``But this is far too complicated.'' \ts{16}{1}{32:33}\ts{16}{1}{32:54}\ts{16}{1}{33:14}\ts{16}{1}{35:25}\ts{16}{1}{35:45}\ts{16}{1}{38:35}\ts{16}{1}{39:08}\ts{16}{1}{39:29}

\begin{definition}[New concepts: optical metric]\label{def:c16-optical}
For a \emph{static} spacetime ($\omega_a=0$) the null condition gives
\[
  \mathrm dt^2=\frac{h_{ab}}{V^2}\,\mathrm dx^a\mathrm dx^b=:\tilde h_{ab}\,\mathrm dx^a\mathrm dx^b,
\]
and $\tilde h_{ab}=h_{ab}/V^2$ is the \textbf{optical metric}: a Riemannian metric on three-dimensional space with signature $(+,+,+)$ (the lecture writes it $\tilde g_{ab}$ or $\tilde h_{ab}$). The \textbf{optical geometry} is the geometry of the spatial light rays: by Fermat's principle \cref{thm:c16-fermat}, \emph{the spatial light rays are the geodesics of the optical metric}. Summary: static spacetimes have a Riemannian optical geometry, stationary ones a Randers--Finsler one. (``Optical'' because it describes light rays; the name comes from Fermat's principle in optics.)
\end{definition}
\ts{16}{1}{33:35}\ts{16}{1}{33:55}\ts{16}{1}{34:15}\ts{16}{1}{34:36}\ts{16}{1}{35:04}\ts{16}{1}{35:45}\ts{16}{1}{36:08}\ts{16}{1}{36:31}\ts{16}{1}{36:52}
In the lecture the line is first written with $\mathrm dt^2$ on the left and $h_{ab}/V^2$ on the right and then named: ``now we have a line element in three-dimensional Riemannian space; this is called the optical metric'' \ts{16}{1}{36:08}\ts{16}{1}{36:31}.

\section{The optical geometry of Schwarzschild}\label{sec:c16-schwarz}

To do calculations take Schwarzschild, ``a particular example of a static spacetime you have met before'' (the rotating black hole would be a stationary one) \ts{16}{1}{38:13}\ts{16}{1}{39:29}. The line element ``in all its glory'' is
\[
  \mathrm ds^2=-\Bigl(1-\frac{2M}{r}\Bigr)\mathrm dt^2+\frac{\mathrm dr^2}{1-2M/r}+r^2\bigl(\mathrm d\theta^2+\sin^2\theta\,\mathrm d\phi^2\bigr),
\]
so $V^2=1-2M/r$ and the spatial metric is $h=\frac{\mathrm dr^2}{1-2M/r}+r^2\mathrm d\Omega^2$. Setting $\mathrm ds^2=0$ and rearranging for $\mathrm dt$ (one looks for null geodesics), the optical metric is
\[
  \tilde h=\frac{\mathrm dr^2}{(1-2M/r)^2}+\frac{r^2}{1-2M/r}\,\mathrm d\Omega^2.
\]
By the spherical symmetry every plane through the centre is a totally geodesic submanifold, so the light rays stay in a plane through the origin; without loss of generality take the equatorial plane $\theta=\pi/2$ (any plane through the origin will do), and the problem is a surface with the diagonal metric
\begin{equation}\label{eq:c16-opt}
  \mathrm dt^2=\frac{\mathrm dr^2}{(1-2M/r)^2}+\frac{r^2\,\mathrm d\phi^2}{1-2M/r}.
\end{equation}
By Fermat's principle the spatial light rays are the geodesics of this metric \ts{16}{1}{39:49}\ts{16}{1}{40:10}\ts{16}{1}{40:31}\ts{16}{1}{40:51}\ts{16}{1}{41:12}\ts{16}{1}{41:33}\ts{16}{1}{41:53}\ts{16}{1}{42:17}.

\paragraph{An isometric embedding.} How can one understand the geodesics of a surface? One way is to \emph{embed it isometrically in $\mathbb R^3$} and look at what it looks like. Use cylindrical coordinates in $\mathbb R^3$ with the Euclidean metric $\mathrm dz^2+\mathrm dR^2+R^2\mathrm d\phi^2$, identify the polar angles $\phi$ of the optical metric and of the embedding, and look for $z$ as a function of $R$. Comparing the coefficients of $\mathrm d\phi^2$ gives the function $R(r)$,
\[
  R=\frac{r}{\sqrt{1-2M/r}},
\]
and comparing the rest, $\frac{\mathrm dr^2}{(1-2M/r)^2}=\bigl(1+z'(R)^2\bigr)\bigl(R'(r)\bigr)^2\mathrm dr^2$ (prime on $R$ is $\mathrm d/\mathrm dr$; ``it is just algebra, you can rearrange it''). One finds
\[
  \Bigl(\frac{\mathrm dz}{\mathrm dR}\Bigr)^2=\frac{4Mr-9M^2}{(r-3M)^2}
\]
[the board reads ``$4Mr-9M^2$ over $(r-3M)^2$'', exact; verified]. This is in terms of the little $r$; one should also express it in terms of the big $R$, but ``that is not easy at all, you really have to integrate, invert and so on''. One does not need to: one can read off what the surface looks like \ts{16}{1}{42:37}\ts{16}{1}{43:01}\ts{16}{1}{43:22}\ts{16}{1}{43:47}\ts{16}{1}{44:07}\ts{16}{1}{44:28}\ts{16}{1}{44:49}\ts{16}{1}{45:09}\ts{16}{1}{45:30}\ts{16}{1}{45:51}.
\begin{enumerate}
  \item \emph{The embedding exists only for $r>9M/4$:} the numerator must be non-negative. Since $9/4=2.25>2$, even though the optical geometry is defined all the way to the event horizon, the embedding \emph{in these coordinates} only works up to $\tfrac94M$, just outside the horizon; physically this is not important \ts{16}{1}{46:12}\ts{16}{1}{46:32}.
  \item \emph{Asymptotics:} far from the black hole $R\approx r$ and $(\mathrm dz/\mathrm dR)^2\to4M/R$, so $z\approx4\sqrt{MR}$: at large distance the surface looks like a square-root function (a paraboloid on its side) \ts{16}{1}{46:52}\ts{16}{1}{47:14}\ts{16}{1}{47:34}\ts{16}{1}{47:56}.
  \item \emph{The throat:} $\mathrm dz/\mathrm dR$ diverges at $r=3M$. In terms of the big $R$ this is not a problem (``you just have to believe me, if you change the coordinates''); what it means is that there is a \textbf{photon orbit}: the radius of a circular photon orbit in Schwarzschild, which one may have met in the tutorials \ts{16}{1}{48:18}\ts{16}{1}{48:39}.
\end{enumerate}
One can then draw it: at large distance $z$ looks like a square root as a function in the plane of $R$ and $\phi$, there is a \textbf{throat} at $r=3M$, and the embedding breaks down a little further in, at $r=\tfrac94M$. ``This is what it looks like: this interesting shape with a throat, a waist here.'' The statement from Fermat's principle is that spatial light rays are the geodesics of this surface: put a light source somewhere on it, an observer somewhere else, and the geodesic between them on the surface is the light ray. Since the embedding is isometric one can talk about the geodesics of the embedded surface and get intuition; the \emph{optical distance} at large distance is just the Euclidean distance (far away in Schwarzschild one is basically in flat space), while closer to the horizon the picture changes \ts{16}{1}{48:39}\ts{16}{1}{49:00}\ts{16}{1}{49:22}\ts{16}{1}{49:44}\ts{16}{1}{50:04}\ts{16}{1}{50:25}\ts{16}{1}{50:48}.
\begin{coursenote}[Student questions on the embedding]
(1) The derivation of $\mathrm dz/\mathrm dR$ is ``a bit tedious, I am not going to go over it'': differentiate $R(r)$, solve for $\mathrm dz/\mathrm dr$ in terms of little $r$, then use the chain rule $\mathrm dz/\mathrm dR=(\mathrm dz/\mathrm dr)/(\mathrm dR/\mathrm dr)$; the real difficulty is to write the right-hand side in terms of big $R$. (2) The region of validity: staticness holds outside $r=2M$; Schwarzschild is perfectly static there, Kerr is only stationary \ts{16}{1}{51:08}\ts{16}{1}{51:28}\ts{16}{1}{51:51}\ts{16}{1}{52:12}\ts{16}{1}{52:33}.
\end{coursenote}

\section{Gaussian curvature}\label{sec:c16-gauss}

``Now we want to discuss lensing, in some sense the deflection of light rays; to do this we have to understand how the curvature of this surface impinges on the propagation of light rays'' \ts{16}{1}{52:54}. The relevant quantity is the \textbf{Gaussian curvature} of the surface: geometrically the product of the two principal curvatures. Looking at the surface above: in one direction it curves in, in the other direction out, so the two principal curvatures have \emph{opposite sign} and the Gaussian curvature is \textbf{negative everywhere}. ``This is very surprising if you think about it; I will come back to this later'' \ts{16}{1}{53:34}\ts{16}{1}{53:54}\ts{16}{1}{54:15}.

In terms of the Riemann tensor (``the notation might be slightly different; you will see that everybody does it differently in books, you just have to get used to it''): for commuting vector fields, $R(X,Y)Z=\nabla_X\nabla_YZ-\nabla_Y\nabla_XZ$, and the components in a basis are read off as $R_{abcd}$ \ts{16}{1}{54:36}\ts{16}{1}{54:57}. On a \emph{two-dimensional} surface, using the symmetries of the Riemann tensor (antisymmetry in each pair, symmetry under the exchange of the pairs) there is just \emph{one} independent component; its value is the Gaussian curvature $K$ (defined up to a sign by convention),
\begin{equation}\label{eq:c16-K}
  R_{abcd}=K\bigl(g_{ad}\,g_{bc}-g_{ac}\,g_{bd}\bigr)
\end{equation}
(``we have to respect the symmetries''; the sign is the one of the lecture and its check is the sphere below) \ts{16}{1}{55:18}\ts{16}{1}{55:39}\ts{16}{1}{56:00}\ts{16}{1}{56:20}.

\begin{lemma}[Gaussian curvature of a surface $\mathrm du^2+f(u)^2\mathrm dv^2$]\label{lem:c16-gauss}
Let a surface be parametrised by coordinates $u,v$ with line element $\mathrm ds^2=\mathrm du^2+f(u)^2\,\mathrm dv^2$ (a surface of revolution; the formula is useful in more general contexts). Then
\[
  K=-\frac{f''(u)}{f(u)}.
\]
\end{lemma}
\ts{16}{1}{56:40}\ts{16}{1}{57:01}\ts{16}{1}{57:23}\ts{16}{1}{57:44}\ts{16}{1}{58:04}
\begin{proof}
One has to work out the Riemann tensor for this very simple metric (``like an exercise''). The metric is diagonal, $\tilde g=\operatorname{diag}(1,f^2)$, with inverse $\operatorname{diag}(1,1/f^2)$. One can convince oneself that there are only two non-vanishing connection components, $\Gamma^u{}_{vv}$ and $\Gamma^v{}_{uv}$ (and $\Gamma^v{}_{vu}$ by symmetry): from $\Gamma^a{}_{bc}=\tfrac12g^{ad}\bigl(g_{db,c}+g_{dc,b}-g_{bc,d}\bigr)$,
\[
  \Gamma^u{}_{vv}=\tfrac12\,g^{uu}\bigl(-g_{vv,u}\bigr)=-\tfrac12\bigl(2ff'\bigr)=-f\,f',
\]
\[
  \Gamma^v{}_{uv}=\tfrac12\,g^{vv}\,g_{vv,u}=\tfrac12\cdot\frac1{f^2}\cdot2ff'=\frac{f'}{f}
\]
(``the off-diagonal terms are zero''; at one point the lecturer first writes the derivative with respect to the wrong variable and corrects it: it is $g_{vv,u}$) \ts{16}{1}{58:27}\ts{16}{1}{58:48}\ts{16}{1}{59:09}\ts{16}{1}{59:31}\ts{16}{1}{59:55}\ts{16}{1}{1:00:17}\ts{16}{1}{1:00:39}\ts{16}{1}{1:01:01}\ts{16}{1}{1:01:22}\ts{16}{1}{1:01:42}.
Now the Riemann tensor. By \eqref{eq:c16-K}, $R_{uvuv}=K\bigl(g_{uv}g_{vu}-g_{uu}g_{vv}\bigr)=-K\,\det\tilde g=-K\,f^2$, so
\[
  K=-\frac{R_{uvuv}}{f^2},
\]
``minus one over $f^2$ times this component''. Pulling the index down with the diagonal metric ($g_{uu}=1$) $R_{uvuv}=R^u{}_{vuv}$, which with the convention chosen in the lecture (``depends on the conventions you choose'': the order of the derivative terms is the opposite of the usual one, which is what makes the result come out right for $K$ as defined above) reads
\[
  R^u{}_{vuv}=\partial_v\Gamma^u{}_{uv}-\partial_u\Gamma^u{}_{vv}+\Gamma^u{}_{ve}\Gamma^e{}_{uv}-\Gamma^u{}_{ue}\Gamma^e{}_{vv}.
\]
The first term vanishes ($\Gamma^u{}_{uv}=0$: there are no connection components with $u$ up and $u$ down together with a $v$); in the product terms with $u$ up and $v$ down the only surviving connection component is the one with $v$ in the middle. So
\[\begin{aligned}
  R^u{}_{vuv}&=-\partial_u(-ff')+\Gamma^u{}_{vv}\Gamma^v{}_{uv}\\
  &=\bigl(f'^2+ff''\bigr)+(-ff')\frac{f'}{f}=f\,f'' ,
\end{aligned}\]
since the $f'^2$ terms cancel. Hence $K=-ff''/f^2=-f''/f$. (The signs in the lecture are first written with an error that is corrected on the board: ``this was minus $ff'$, yes, that makes more sense''; the final cancellation is the one above: ``we get plus $ff''$ plus $f'^2$ minus $f'^2$, these cancel'') \ts{16}{1}{1:02:06}\ts{16}{1}{1:02:27}\ts{16}{1}{1:02:48}\ts{16}{1}{1:03:08}\ts{16}{1}{1:03:28}\ts{16}{1}{1:03:50}\ts{16}{1}{1:04:11}\ts{16}{1}{1:04:31}\ts{16}{1}{1:04:54}\ts{16}{1}{1:05:17}\ts{16}{1}{1:05:38}.
\end{proof}

\begin{example}[The round sphere as a check]\label{ex:c16-sphere}
The unit $S^2$ has $\mathrm ds^2=\mathrm d\vartheta^2+\sin^2\vartheta\,\mathrm d\varphi^2$, i.e.\ $u=\vartheta$, $v=\varphi$, $f(u)=\sin u$. The Gaussian curvature of a sphere is one over the radius of curvature squared, since the \emph{osculating sphere} of a surface (``osculating, from the Latin \emph{osculari}, to kiss'' [etymology as said in class, not checked]: it touches the surface not just in the first derivative but also in the second) is the sphere itself; so one expects $K=+1$. By the lemma $K=-f''/f=-(-\sin u)/\sin u=+1$ \ts{16}{1}{1:05:59}\ts{16}{1}{1:06:20}\ts{16}{1}{1:06:40}\ts{16}{1}{1:07:01}\ts{16}{1}{1:07:23}.
\end{example}

\paragraph{Application to the Schwarzschild optical geometry.} After the break one does the lensing, but first one applies the lemma to \eqref{eq:c16-opt}: identify $v=\phi$ (the polar coordinate again), and $\mathrm du=\mathrm dr/(1-2M/r)$, $f(r)=r/\sqrt{1-2M/r}$ (so that $f^2=r^2/(1-2M/r)$ is read off from the line element). Expressing $f$ as a function of $u$ would mean integrating, changing coordinates, differentiating: ``tedious, let us be a bit more clever''. Use $K=-\frac1f\frac{\mathrm d^2f}{\mathrm du^2}$ and the chain rule, $\frac{\mathrm d}{\mathrm du}=\frac{\mathrm dr}{\mathrm du}\frac{\mathrm d}{\mathrm dr}$ with $\mathrm dr/\mathrm du=1-2M/r$:
\[
  \frac{\mathrm d^2f}{\mathrm du^2}=\frac{\mathrm dr}{\mathrm du}\frac{\mathrm d}{\mathrm dr}\Bigl(\frac{\mathrm dr}{\mathrm du}\frac{\mathrm df}{\mathrm dr}\Bigr)=\Bigl(\frac{\mathrm dr}{\mathrm du}\Bigr)\frac{\mathrm d}{\mathrm dr}\Bigl(\frac{\mathrm dr}{\mathrm du}\Bigr)\frac{\mathrm df}{\mathrm dr}+\Bigl(\frac{\mathrm dr}{\mathrm du}\Bigr)^2\frac{\mathrm d^2f}{\mathrm dr^2},
\]
``differentiating is always easier than integrating''. Plugging in $f(r)=r/\sqrt{1-2M/r}$ and working it out, ``after some work'':
\begin{equation}\label{eq:c16-KS}
  K=-\frac{2M}{r^3}\Bigl(1-\frac{3M}{2r}\Bigr).
\end{equation}
(Checked symbolically with sympy.) The bracket is always positive outside $r=2M$ (since $3M/2r<3/4$ there), so \emph{the Gaussian curvature of the Schwarzschild optical geometry is negative everywhere}: as already clear from the principal curvatures in the picture, one can now also see it analytically, ``going through all these shenanigans'' \ts{16}{1}{1:07:43}\ts{16}{1}{1:08:06}\ts{16}{1}{1:08:27}\ts{16}{1}{1:08:48}\ts{16}{1}{1:09:09}\ts{16}{1}{1:09:30}\ts{16}{1}{1:09:50}\ts{16}{1}{1:10:15}\ts{16}{1}{1:10:36}\ts{16}{1}{1:10:56}\ts{16}{1}{1:11:17}\ts{16}{1}{1:11:37}\ts{16}{1}{1:11:58}\ts{16}{1}{1:12:19}\ts{16}{1}{1:12:41}.
\begin{coursenote}[The location of the zero of the bracket]
The lecture says ``this is three-halves $M$'' for the point where the bracket would vanish ($r=\tfrac32M$): this lies inside the horizon $2M$, so outside $2M$ the bracket never vanishes \ts{16}{1}{11:58}.
\end{coursenote}

\paragraph{The question.} Before the break one observation: the Gaussian curvature is negative, and \emph{the geodesics on the surface are our light rays}; negative Gaussian curvature means (to be shown in the next lecture) that the geodesics \emph{diverge}, and they diverge everywhere. How is it then possible for light rays to reconverge somewhere else, so that one sees multiple images as in the pictures of yesterday? The answer is that this is a \emph{global effect}: one has to connect the local property of the surface, the curvature, with the global property, the topology, and this is done by the Gauss--Bonnet theorem, ``after the break'' (chapter~\ref{ch:c17}) \ts{16}{1}{1:13:08}\ts{16}{1}{1:13:32}\ts{16}{1}{1:13:53}.

\section*{Key concepts}
\begin{keyconcepts}
\textbf{Retrieval questions.}
\begin{enumerate}
  \item State Fermat's principle in classical optics. Why does it not work for null curves in general relativity?
  \item Define stationary and static spacetimes. Write the general stationary line element and say what $\omega_a$ is.
  \item Derive \eqref{eq:c16-deltat} and explain why the arrival time, not the total travel time, is stationary. Why does staticness matter?
  \item What is the optical metric? Which geometry do stationary spacetimes have?
  \item Compute the optical metric of Schwarzschild and its embedding in $\mathbb R^3$: where does it break down, and what is special at $r=3M$?
  \item State and prove $K=-f''/f$, and check it on the sphere. Why is the Gaussian curvature of the Schwarzschild optical geometry negative everywhere?
\end{enumerate}
\textbf{Concepts.}
\begin{itemize}
  \item Stationary and static spacetimes, null geodesics (\cref{def:c16-stat}, \cref{def:c16-static}).
  \item Fermat's principle in static spacetimes: stationary arrival time (\cref{thm:c16-fermat}).
  \item Optical metric and optical geometry; Randers--Finsler for stationary spacetimes (\cref{def:c16-optical}).
  \item Gaussian curvature and its formula $-f''/f$ (\cref{lem:c16-gauss}); the Schwarzschild value \eqref{eq:c16-KS}.
\end{itemize}
\end{keyconcepts}

\flushside
\end{paracol}
